%% ============================================================
%% A Curvature-Scale Coincidence on Compact S^3
%% and Two Dynamical Completion Tests
%% Technical preprint v1.3.5-final-checked
%% Andrei Preece & Boris Batenin -- Berlin, 2026
%% ============================================================
\documentclass[11pt,a4paper]{article}

\usepackage[T1]{fontenc}
\usepackage[utf8]{inputenc}
\usepackage{lmodern}
\usepackage{microtype}
\usepackage{amsmath,amssymb,amsfonts,mathtools}
\usepackage{bm}
\usepackage{graphicx}
\usepackage{booktabs}
\usepackage{array}
\usepackage{longtable}
\usepackage{siunitx}
\usepackage{hyperref}
\usepackage{geometry}
\usepackage{xcolor}
\usepackage{caption}
\usepackage{subcaption}
\usepackage{enumitem}
\usepackage{tcolorbox}
\usepackage{fancyhdr}
\usepackage{titlesec}
\usepackage{parskip}
\usepackage[numbers,sort&compress]{natbib}

\geometry{margin=1in, top=2.2cm, bottom=2.2cm}
\setlength{\headheight}{15pt}
\sisetup{per-mode=symbol, detect-all=true, exponent-product=\times}
\tcbuselibrary{most}

\definecolor{resultcol}{RGB}{0,90,160}
\definecolor{caveatcol}{RGB}{170,60,0}
\definecolor{openpcol}{RGB}{130,0,130}
\definecolor{falsecol}{RGB}{0,120,50}
\definecolor{defcol}{RGB}{50,50,100}

\newtcolorbox{resultbox}[1][Result]{
  colback=blue!5,colframe=resultcol,fonttitle=\bfseries\small,
  title=#1,boxrule=0.7pt,arc=3pt,left=4pt,right=4pt,top=4pt,bottom=4pt}
\newtcolorbox{caveatbox}[1][Caveat]{
  colback=orange!7,colframe=caveatcol,fonttitle=\bfseries\small,
  title=#1,boxrule=0.7pt,arc=3pt,left=4pt,right=4pt,top=4pt,bottom=4pt}
\newtcolorbox{openpbox}[1][Open Problem]{
  colback=purple!5,colframe=openpcol,fonttitle=\bfseries\small,
  title=#1,boxrule=0.7pt,arc=3pt,left=4pt,right=4pt,top=4pt,bottom=4pt}
\newtcolorbox{falsebox}[1][Falsifiability Test]{
  colback=green!5,colframe=falsecol,fonttitle=\bfseries\small,
  title=#1,boxrule=0.7pt,arc=3pt,left=4pt,right=4pt,top=4pt,bottom=4pt}
\newtcolorbox{defbox}[1][Definition]{
  colback=gray!5,colframe=defcol,fonttitle=\bfseries\small,
  title=#1,boxrule=0.6pt,arc=3pt,left=4pt,right=4pt,top=4pt,bottom=4pt}

\hypersetup{colorlinks=true,linkcolor=resultcol,citecolor=resultcol,urlcolor=resultcol}

\pagestyle{fancy}
\fancyhf{}
\fancyhead[L]{\small\textcolor{resultcol}{Curvature-Scale Coincidence on Compact $S^3$}}
\fancyhead[R]{\small\textcolor{gray}{Preece \& Batenin, 2026 -- v1.3.5}}
\fancyfoot[C]{\thepage}
\renewcommand{\headrulewidth}{0.4pt}

\newcommand{\vph}{\varphi}
\newcommand{\ellP}{\ell_{\rm P}}
\newcommand{\rhoP}{\rho_{\rm P}}
\newcommand{\rhoL}{\rho_{\Lambda,{\rm obs}}}
\newcommand{\rhoC}{\rho_{\rm crit}}
\newcommand{\RH}{R_H}
\newcommand{\Reff}{R_{\rm eff}}
\newcommand{\Rcurv}{R_{\rm curv}}
\newcommand{\RS}{R_{S^3}}
\newcommand{\OmL}{\Omega_\Lambda}
\newcommand{\OmK}{\Omega_k}
\newcommand{\CEH}{C_{\rm EH}}
\newcommand{\Cobs}{C_{\rm obs}}
\newcommand{\Mpc}{\mathrm{Mpc}}
\newcommand{\Gly}{\mathrm{Gly}}
\newcommand{\sech}{\operatorname{sech}}

\begin{document}

\begin{titlepage}
\begin{center}
\vspace*{0.7cm}
{\fontsize{18}{22}\bfseries\color{resultcol}
  A Curvature-Scale Coincidence on Compact $S^3$}\\[0.5em]
{\fontsize{15}{19}\bfseries
  and Two Dynamical Completion Tests}

\vspace{0.6cm}
{\large\itshape\color{gray}
  Technical preprint on a constant-positive-curvature $R^{-2}$ scale}

\vspace{0.8cm}
\rule{0.85\textwidth}{1.2pt}

\vspace{0.4cm}
{\large Andrei Preece\textsuperscript{1} \quad Boris Batenin\textsuperscript{1}}\\[0.3em]
{\normalsize\textsuperscript{1}Independent Researchers, Berlin, Germany}\\[0.15em]
{\small\texttt{andre.preece@web.de} \quad \texttt{batenin@web.de}}

\vspace{0.4cm}
\rule{0.85\textwidth}{1.2pt}

\vspace{0.8cm}
\begin{minipage}{0.88\textwidth}
\textbf{Abstract.}
We study compact, simply connected spatial geometry $\Sigma\simeq S^3$ and the
associated curvature-density scale on $S^3(R)$.  The algebraic core is the
identity $\rho_P\ell_P^2=c^4/G$: in the tested integer family
$\rho_\beta(n)=\rho_P\varphi^{-\beta n}$, the branch $\beta=2$ is exactly
$\rho_2(R)=c^4/(G R^2)$ and is the only branch among $\beta=1,2,3,4$ that reaches
the observed dark-energy density at cosmological radii.  The $\varphi$ ladder is
therefore used here as a logarithmic bookkeeping coordinate unless an independent
integer-level selection principle is supplied.  For constant positive curvature, and in particular on compact $S^3(R)$, the spatial
Einstein--Hilbert scalar-curvature integral satisfies
$\int_{S^3}\sqrt{g}\,{}^{(3)}R\,d^3x=12\pi^2R$, giving the natural positive
scalar-curvature scale
\[\rho_{\rm EH}^{(3)}(R)=\frac{3}{8\pi}\frac{c^4}{G R^2}.\]
The coefficient $3/(8\pi)$ is a constant-positive-curvature density coefficient;
compactness makes the total integral finite but does not make the local density
coefficient topological.  Evaluated at the observationally defined radius
$\Reff=\RH/\sqrt{\OmL}=c/(H_0\sqrt{\OmL})$, this reproduces
$\rho_\Lambda=\OmL\rho_c$ exactly because $\Reff$ is defined from the observed
dark-energy scale.  This equality is not a derivation of dark energy: $\OmL$ and
$\Reff$ are inserted from the Friedmann decomposition, and a conserved $R^{-2}$
component has $w=-1/3$, not $w\simeq-1$.  We therefore separate four layers: (i) the robust algebraic scaling
result, (ii) the geometrically normalized positive scalar-curvature scale,
(iii) the observationally inserted radius identification, and (iv) two incomplete
dynamical completion tests--a CMC/York phase-drift ansatz and post-recombination
null-domain kinematics, whose product-rule identity is
$\ddot D_\gamma=\ddot R\chi_\gamma+c\dot R/R$.  The resulting program is
falsifiable through global
curvature bounds, closed-geometry distances, low-$\ell$ CMB/topology tests,
Pantheon+ supernovae, DESI DR2 BAO, and the sign and magnitude of the
dark-energy equation-of-state evolution.
\end{minipage}

\vspace{0.7cm}
\textbf{Keywords:} compact $S^3$ topology; vacuum energy; curvature density;
$\vph$-bookkeeping; Einstein--Hilbert action; CMC phase law; null-domain kinematics;
dark energy

\vspace{0.5cm}
\textit{Preprint v1.3.5 -- May 2026.}
\end{center}
\end{titlepage}

\tableofcontents
\newpage

\section{Introduction and Scope}
\label{sec:intro}

The observed dark-energy density is of order
$\rhoL\simeq5.25\times10^{-10}\,\mathrm{J\,m^{-3}}$ for the Planck 2018
baseline values $H_0\simeq67.4\,\mathrm{km\,s^{-1}\,Mpc^{-1}}$ and
$\OmL\simeq0.685$~\cite{Planck2018}.  Its smallness relative to Planck-scale
vacuum estimates is the cosmological-constant problem~\cite{Wald1984}.
The aim of this note is narrow: to test whether compact positive curvature and
the associated $R^{-2}$ scaling can reproduce the \emph{scale} of $\rhoL$ without
claiming a complete derivation of late-time acceleration.

The proposed geometry is a closed spatial section $\Sigma\simeq S^3$.
The proposed scale ladder is
\begin{equation}
R_n=\ellP\vph^n,\qquad \vph=\frac{1+\sqrt5}{2}.
\end{equation}
The main result is that the branch $\rho\propto R^{-2}$ is uniquely selected
within the small integer family tested in Sec.~\ref{sec:beta}.  The same branch
is generated by the spatial Einstein--Hilbert scalar-curvature integral on $S^3$.
This establishes a curvature-scale coincidence, not a derivation of the
late-time dark-energy component.

\begin{defbox}[Status convention]
The paper uses three labels.  \emph{Result} means an algebraic, numerical, or
geometric statement derived in the paper.  \emph{Candidate} means a proposed
completion requiring additional derivation or data fitting.  \emph{Open problem}
means no claim of solution is made here.
\end{defbox}

\textbf{Nonclaims.}  We do not prove that dark energy is topology, do not remove
$\Lambda$ from the gravitational action, do not infer cosmic topology from
large-scale structure alone, and do not identify $\Reff$ with the global FLRW
curvature radius.  The result is a geometrically normalized scale-setting observation plus a set of
conservative dynamical tests.

\section{Compact \texorpdfstring{$S^3$}{S3} Geometry}
\label{sec:geom}

\begin{defbox}[Spatial topology]
The spatial manifold is assumed to be compact, orientable, boundaryless and
simply connected: $\Sigma\simeq S^3$.  By the Poincare--Perelman theorem, a
closed simply connected 3-manifold is homeomorphic to $S^3$~\cite{Perelman2002,Perelman2003}.
Spherical quotient spaces $S^3/\Gamma$ remain possible extensions, but the
minimal case studied here is the simply connected one.
\end{defbox}

With curvature radius $R(t)$, the closed FLRW line element is
\begin{equation}
 ds^2=-c^2dt^2+R^2(t)\left[d\chi^2+\sin^2\chi\,d\Omega_2^2\right],
 \qquad 0\leq\chi\leq\pi .
\label{eq:metric}
\end{equation}
When comparison with FLRW observables is required, we write
$R(t)=a(t)R_c$ with $a_0=1$, so that $R_c$ is the present global curvature
radius.  The effective normalization radius $\Reff$ introduced below is not
identified with $R_c$.

The standard geometric quantities are
\begin{align}
V_{S^3}(R)&=2\pi^2R^3,\label{eq:VS3}\\
{}^{(3)}R&=\frac{6}{R^2},\label{eq:3R}\\
V_{\rm ball}(\chi,R)&=2\pi R^3\left(\chi-\frac12\sin2\chi\right),\\
A(\chi,R)&=4\pi R^2\sin^2\chi .
\end{align}
The dimensionless volume-growth kernel of a geodesic ball is
\begin{equation}
K(\chi)=\frac{1}{V_{\rm ball}}\frac{dV_{\rm ball}}{d\chi}
=\frac{2\sin^2\chi}{\chi-\sin\chi\cos\chi}.
\label{eq:kernel}
\end{equation}
It has $K(\chi)\sim3/\chi$ as $\chi\to0$ and $K(\chi)\to0$ as $\chi\to\pi$.

\section{The \texorpdfstring{$\varphi$}{phi}-Hierarchy and the Two-Index Correction}
\label{sec:hier}

The geometric hierarchy is
\begin{equation}
R_n=\ellP\vph^n,
\qquad
V_n=2\pi^2\ellP^3\vph^{3n},
\qquad
A_n=4\pi\ellP^2\vph^{2n},
\qquad
n_R(R)=\frac{\ln(R/\ellP)}{\ln\vph}.
\label{eq:hier}
\end{equation}
\begin{caveatbox}[$\varphi$-hierarchy status]
Because $n_R=\ln(R/\ell_P)/\ln\varphi$ is allowed to be non-integer in the present
paper, the golden-ratio hierarchy is currently a logarithmic coordinate and
bookkeeping device.  It does not by itself explain the vacuum scale.  A physical
$\varphi$-hierarchy would require an independent principle selecting integer
levels, phase locking, spectral gaps, or a genuinely $\varphi$-recursive
dynamical structure.
\end{caveatbox}
For the Planck 2018 baseline,
\begin{equation}
 n_R(\RH)=291.5436,
 \qquad
 n_R(\Reff)=291.9367.
\end{equation}
Thus a cosmic radius is a geometric level near $292$, not near $85$.

\begin{caveatbox}[Level 85 is not a geometric cosmic radius]
Under $R_n=\ellP\vph^n$, one finds
\begin{equation}
R_{85}=9.3855\times10^{-18}\,\mathrm{m},
\end{equation}
which is sub-nuclear.  Any separate ``85-level'' used in auxiliary notes must
therefore be interpreted as a distinct amplitude or calibration index, not as
$R_{85}$ in the geometric ladder.
\end{caveatbox}

A possible two-index relation is
\begin{equation}
 n_A=n_R-n_{\rm cal},\qquad n_{\rm cal}\simeq207,
\label{eq:ncal}
\end{equation}
for which $R_{207}=2.9441\times10^8\,\mathrm{m}$ and
$n_A(\RH)\simeq84.54$.  Equation~\eqref{eq:ncal} is only a calibration
hypothesis.

\section{Density Scaling Family and Selection of \texorpdfstring{$\beta=2$}{beta=2}}
\label{sec:beta}

Let $E_n=E_P\vph^{\alpha n}$ and divide by $V_n\propto\vph^{3n}$.  Then
\begin{equation}
\rho_n=\rhoP\vph^{-\beta n},\qquad \beta=3-\alpha,
\label{eq:rho_beta}
\end{equation}
where the Planck energy density is
\begin{equation}
\rhoP=\frac{c^7}{\hbar G^2}=4.6329\times10^{113}\,\mathrm{J\,m^{-3}}.
\end{equation}
The matching level and radius for a target density $\rhoL$ are
\begin{equation}
n^*_\beta=\frac{\ln(\rhoP/\rhoL)}{\beta\ln\vph},
\qquad
R^*_\beta=\ellP\vph^{n^*_\beta}.
\label{eq:nstar}
\end{equation}

\begin{table}[htbp]
\centering
\caption{$\beta$-branch matching table.  The radius labels are corrected: $\beta=3$ is Earth-scale, and $\beta=4$ is sub-millimetre.}
\label{tab:beta}
\small
\begin{tabular}{@{}cllrrl@{}}
\toprule
$\beta$ & $\alpha$ & $\rho\propto$ & $n^*_{\beta}$ & $R^*_{\beta}$ & Status \\
\midrule
1 & 2 & $R^{-1}$ & 588.2904 & $1.43\times10^{88}\,\mathrm{m}$ & excluded: super-cosmic \\
\textbf{2} & \textbf{1} & $\mathbf{R^{-2}}$ & \textbf{294.1452} & $\mathbf{4.80\times10^{26}\,\mathrm{m}=50.73\,Gly}$ & \textbf{scale-compatible} \\
3 & 0 & $R^{-3}$ & 196.0968 & $1.55\times10^6\,\mathrm{m}$ & excluded: planetary \\
4 & $-1$ & $R^{-4}$ & 147.0726 & $8.81\times10^{-5}\,\mathrm{m}$ & excluded: sub-mm \\
\bottomrule
\end{tabular}
\end{table}

Because $\rhoP\ellP^2=c^4/G$, the $\beta=2$ branch is exactly
\begin{equation}
\rho_2(R)=\frac{c^4}{GR^2}.
\label{eq:C1scale}
\end{equation}
The coefficient in Eq.~\eqref{eq:C1scale} is unity.  The spatial
Einstein--Hilbert normalization derived below supplies a geometric coefficient
$\CEH=3/(8\pi)$, moving the equality radius from $50.73\,\Gly$ to
$17.53\,\Gly$.

\begin{resultbox}[Scaling-level result]
Within the tested integer family $\beta=1,2,3,4$, only $\beta=2$ places the
observed vacuum-density scale at cosmological radii.  This is a scaling result,
not a derivation of the dark-energy equation of state.
\end{resultbox}

\begin{figure}[htbp]
\centering
\includegraphics[width=0.92\textwidth]{figures/fig1_density_scaling.pdf}
\caption{Density scaling $\rho_\beta(n)=\rhoP\vph^{-\beta n}$ for
$\beta=1,2,3,4$.  Only $\beta=2$ intersects $\rhoL$ at a cosmological level.}
\label{fig:density}
\end{figure}

\section{Einstein--Hilbert Curvature Normalization on \texorpdfstring{$S^3$}{S3}}
\label{sec:EH}

On $S^3(R)$,
\begin{equation}
\int_{S^3}\sqrt{g}\,{}^{(3)}R\,d^3x
=\frac{6}{R^2}\,2\pi^2R^3=12\pi^2R.
\label{eq:integral}
\end{equation}
The integral scales linearly with $R$.  Thus an EH-type curvature-energy scale
is proportional to $R$, and the corresponding density is proportional to
$R^{-2}$.

Define
\begin{equation}
E_{\rm curv}^{(3)}(R)
=\frac{c^4}{16\pi G}\int_{S^3}\sqrt{g}\,{}^{(3)}R\,d^3x
=\frac{3\pi c^4}{4G}R.
\label{eq:Ecurv}
\end{equation}
Dividing by $V_{S^3}=2\pi^2R^3$ gives
\begin{equation}
\rho_{\rm EH}^{(3)}(R)
=\frac{E_{\rm curv}^{(3)}}{V_{S^3}}
=\frac{3}{8\pi}\frac{c^4}{GR^2}.
\label{eq:rho_EH}
\end{equation}

\begin{resultbox}[Geometric normalization]
\begin{equation}
\boxed{\CEH=\frac{3}{8\pi}=0.1193662073.}
\end{equation}
This is the constant-positive-curvature scalar-density coefficient.  Compactness
makes the total spatial integral finite; it does not by itself make the local
density coefficient topological.  The same density coefficient is obtained on
spherical quotients $S^3/\Gamma$ after division by their volume.
\end{resultbox}

\begin{caveatbox}[Hamiltonian status]
Equation~\eqref{eq:rho_EH} is a spatial EH curvature-density scale.  It is not a
complete ADM Hamiltonian energy, since the full Hamiltonian constraint also
contains extrinsic curvature, lapse dependence, matter, and possible subtraction
or reference terms.  The paper therefore uses $\rho_{\rm EH}^{(3)}$ as a
scale-setting density, not as a closed gravitational energy theorem.  In particular,
no conserved quasi-local or global gravitational energy is being claimed for the
compact slice.
\end{caveatbox}

\begin{caveatbox}[Sign of the FLRW curvature fluid]
The positive sign of $\rho_{\rm EH}^{(3)}$ is a convention for a scalar-curvature
scale.  It is not the sign of the effective FLRW curvature density.  In the
standard closed Friedmann equation,
\begin{equation}
H^2=\frac{8\pi G}{3c^2}\rho-\frac{c^2}{a^2R_c^2},
\end{equation}
so moving the curvature term to the matter side gives
\begin{equation}
\rho_k=-\frac{3c^4}{8\pi G a^2R_c^2}\quad(k=+1),
\end{equation}
which is negative.  Therefore $\rho_{\rm EH}^{(3)}$ must not be confused with the
standard closed-curvature fluid.
\end{caveatbox}

\section{Critical-Density Identity and the Tautology Warning}
\label{sec:tautology}

The Hubble radius is $\RH=c/H_0$.  With energy density units,
\begin{equation}
\rhoC=\frac{3H_0^2c^2}{8\pi G}
      =\frac{3}{8\pi}\frac{c^4}{G\RH^2}
      =\CEH\frac{c^4}{G\RH^2}.
\label{eq:crit_EH}
\end{equation}
The observed dark-energy density is
\begin{equation}
\rhoL=\OmL\rhoC
=\frac{3\OmL}{8\pi}\frac{c^4}{G\RH^2}
=\Cobs\frac{c^4}{G\RH^2},
\qquad
\frac{\Cobs}{\CEH}=\OmL.
\label{eq:rhoL}
\end{equation}

\begin{caveatbox}[Tautology warning]
Substituting $R=\RH=c/H_0$ into $c^4/(GR^2)$ reproduces the critical density by
definition.  Likewise, $\rho_{\rm EH}^{(3)}(\Reff)=\rho_\Lambda$ is exact only
because $\Reff=\RH/\sqrt{\OmL}=c/(H_0\sqrt{\OmL})$ is the observed
dark-energy length written in this normalization.  Equivalently, with
$H_\Lambda=H_0\sqrt{\OmL}$, one has $\Reff=c/H_\Lambda$.  The construction
therefore reproduces $\rho_\Lambda$ because $\Reff$ is defined from
$\rho_\Lambda$.  The nontrivial task is to derive $\OmL$, $\Reff$, or a
dynamical selection rule from compact geometry rather than inserting them from
data.
\end{caveatbox}

Define the effective curvature-energy radius by
\begin{equation}
\rho_{\rm EH}^{(3)}(\Reff)=\rhoL.
\end{equation}
Using Eq.~\eqref{eq:rhoL},
\begin{equation}
\boxed{\Reff=\frac{\RH}{\sqrt{\OmL}}=1.6583\times10^{26}\,\mathrm{m}=17.53\,\Gly.}
\label{eq:Reff}
\end{equation}
This equality is exact given the definitions and the observed input $\OmL$; it is the observed dark-energy length written in the present normalization.
It is not a first-principles derivation of $\OmL$.

\begin{figure}[htbp]
\centering
\includegraphics[width=0.88\textwidth]{figures/fig2_rho_EH.pdf}
\caption{$\rho_{\rm EH}^{(3)}(R)=\CEH c^4/(GR^2)$.  The equality with
$\rhoL$ occurs at $\Reff=\RH/\sqrt{\OmL}$; the equality with $\rhoC$ at $\RH$ is
the critical-density identity.}
\label{fig:rho_EH}
\end{figure}

\section{Three Radii That Must Not Be Conflated}
\label{sec:radii}

The construction contains three different radii.

\textbf{Hubble radius.}
\begin{equation}
\RH=\frac{c}{H_0}=1.3725\times10^{26}\,\mathrm{m}=14.51\,\Gly .
\end{equation}

\textbf{Effective curvature-energy radius.}
This is $\Reff=\RH/\sqrt{\OmL}=17.53\,\Gly$.  If the Hubble radius is represented
as an arc length on an auxiliary $S^3(\Reff)$, the angle is
\begin{equation}
\chi_H=\frac{\RH}{\Reff}=\sqrt{\OmL}=0.827647.
\label{eq:chiH}
\end{equation}
The corresponding geodesic-ball volume fraction is
\begin{equation}
f_{\rm obs}(\chi_H)=\frac{V_{\rm ball}(\chi_H,\Reff)}{V_{S^3}(\Reff)}
=\frac{\chi_H-\frac12\sin2\chi_H}{\pi}=0.104861.
\label{eq:fobs}
\end{equation}
This fraction is a diagnostic number, not yet a derived $\OmL$.

\textbf{Global FLRW curvature radius.}
For a closed FLRW model with present scale factor $a_0=1$,
\begin{equation}
\Rcurv=\frac{\RH}{\sqrt{|\OmK|}}.
\label{eq:Rcurv}
\end{equation}
Planck combined with BAO is consistent with spatial flatness, with constraints
of order $|\OmK|\lesssim10^{-3}$--$10^{-2}$ depending on data and model
choices~\cite{Planck2018}.  Thus the global $S^3$ radius must be much larger
than $\RH$.

\begin{caveatbox}[Curvature conflict if $\Reff$ is misidentified]
If $\Reff$ were identified with the global FLRW curvature radius, then
\begin{equation}
\OmK=-\left(\frac{\RH}{\Reff}\right)^2=-\OmL\simeq-0.685,
\end{equation}
which is excluded.  Therefore $\Reff$ can only be an effective curvature-energy
scale, not the global FLRW curvature radius.
\end{caveatbox}

\begin{table}[htbp]
\centering
\caption{Physically distinct radii and $\vph$ levels.}
\label{tab:radii}
\small
\begin{tabular}{@{}lrrrrl@{}}
\toprule
Description & $R$ [m] & $R$ [Gly] & $n_R$ & $n_R-207$ & Role \\
\midrule
$\RH=c/H_0$ & $1.3725\times10^{26}$ & 14.51 & 291.5436 & 84.5436 & observational \\
$\Reff=\RH/\sqrt{\OmL}$ & $1.6583\times10^{26}$ & 17.53 & 291.9367 & 84.9367 & effective scale \\
$\Rcurv(|\OmK|=10^{-2})$ & $1.3725\times10^{27}$ & 145.1 & 296.3285 & 89.3285 & global lower bound \\
$\Rcurv(|\OmK|=10^{-3})$ & $4.3402\times10^{27}$ & 458.8 & 298.7210 & 91.7210 & global lower bound \\
$\Rcurv(|\OmK|=10^{-4})$ & $1.3725\times10^{28}$ & 1450.7 & 301.1135 & 94.1135 & global lower bound \\
\bottomrule
\end{tabular}
\end{table}

\begin{figure}[htbp]
\centering
\includegraphics[width=0.88\textwidth]{figures/fig3_radii.pdf}
\caption{The Hubble radius, the effective curvature-energy radius, and two
illustrative global curvature radii.}
\label{fig:radii}
\end{figure}

\section{The Equation-of-State Gap}
\label{sec:eosgap}

If $\rho_{\rm eff}(a)=\rho_0a^{-2}$ is treated as a conserved perfect-fluid
energy density, the continuity equation gives
\begin{equation}
\rho(a)\propto a^{-3(1+w)},
\qquad
-3(1+w)=-2,
\qquad
w=-\frac13.
\label{eq:w13}
\end{equation}
The acceleration equation in energy-density units is
\begin{equation}
\frac{\ddot a}{a}=-\frac{4\pi G}{3c^2}\rho(1+3w).
\label{eq:ray}
\end{equation}
For $w=-1/3$ this gives no acceleration from that component alone.

\begin{resultbox}[Equation-of-state gap]
The branch $\rho\sim R^{-2}$ gives the correct magnitude near cosmic radii, but
as a conserved fluid it has $w=-1/3$, not $w\simeq-1$.  A second dynamical
principle is required if the branch is to model late-time acceleration.
\end{resultbox}

The conservative interpretation is therefore not
``dark energy is simply $a^{-2}$ curvature fluid.''  Instead, the present paper
uses $\rho_{\rm EH}^{(3)}(\Reff)$ as a normalization scale and asks whether a
separate dynamical factor can keep the effective component nearly constant.

\section{Candidate Test I: Post-Recombination Null-Domain Kinematics}
\label{sec:nulldom}

After recombination, photons propagate freely and define an observable null
domain.  A radial null geodesic of Eq.~\eqref{eq:metric} satisfies
\begin{equation}
\dot\chi_\gamma=\frac{c}{R(t)}.
\label{eq:chidot}
\end{equation}
Define a light-front arc length
\begin{equation}
D_\gamma(t)=R(t)\chi_\gamma(t).
\label{eq:Dgamma}
\end{equation}
The product rule gives
\begin{equation}
\dot D_\gamma=\dot R\chi_\gamma+R\dot\chi_\gamma
=\dot R\chi_\gamma+c,
\label{eq:Dgamma_dot}
\end{equation}
and
\begin{equation}
\ddot D_\gamma
=\ddot R\chi_\gamma+2\dot R\dot\chi_\gamma+R\ddot\chi_\gamma.
\end{equation}
Since $\dot\chi_\gamma=c/R$,
\begin{equation}
\ddot\chi_\gamma=-\frac{c\dot R}{R^2},
\end{equation}
so
\begin{equation}
\boxed{\ddot D_\gamma=\ddot R\chi_\gamma+\frac{c\dot R}{R}.}
\label{eq:Dgamma_ddot}
\end{equation}
For a coasting scale $\ddot R=0$ with $\dot R>0$, this coordinate light-front
length has $\ddot D_\gamma=c\dot R/R>0$.

\begin{caveatbox}[Interpretation limit]
Equation~\eqref{eq:Dgamma_ddot} is a kinematic identity for the chosen physical
arc-length variable, and by itself it does not distinguish the model from a
standard FLRW description.  Until $D_\gamma$ is mapped into $D_C,D_M,D_L,D_A$
and fitted against redshift observables, it is a coordinate-domain identity
rather than a falsifiability test.
\end{caveatbox}

The observable null-domain volume is
\begin{equation}
V_\gamma(t)=2\pi R^3(t)\left[\chi_\gamma(t)-\frac12\sin2\chi_\gamma(t)\right],
\qquad
\frac{1}{R}\,K(\chi_\gamma)=\frac{1}{V_\gamma}\frac{\partial V_\gamma}{\partial \chi_\gamma}.
\end{equation}
Figure~\ref{fig:kernel_null} shows the kernel and toy light-front behavior.

\begin{figure}[htbp]
\centering
\begin{subfigure}[b]{0.48\textwidth}
\includegraphics[width=\textwidth]{figures/fig4_kernel.pdf}
\caption{$S^3$ volume-growth kernel.}
\end{subfigure}
\hfill
\begin{subfigure}[b]{0.48\textwidth}
\includegraphics[width=\textwidth]{figures/fig5_null_kinematics.pdf}
\caption{Toy light-front growth.}
\end{subfigure}
\caption{Null-domain diagnostics.  The plotted light-front curves are illustrative
only and are not fits to data.}
\label{fig:kernel_null}
\end{figure}

\section{Candidate Test II: CMC/York Phase Drift}
\label{sec:cmc}

The second completion is a phase-clock ansatz in a compact CMC/York treatment.
York-time quantum cosmology has been studied in related canonical contexts
\cite{RoserValentini2014}.  The present paper uses only the following ansatz as a candidate completion:
\begin{equation}
H(t)=\frac{1}{3}\dot{\Theta}(t)+\mathcal{O}(\epsilon_{\rm BO}^2),
\label{eq:phaselaw}
\end{equation}
where $\Theta$ is a gauge-invariant determinant-line geometric phase for a
transported transverse-traceless sector.  Equation~\eqref{eq:phaselaw} is not
assumed to be established by existing quantum gravity; it is a candidate closure condition requiring independent derivation.

If $\dot{\Theta}\to\Theta_1$ at late times, then $H\to H_\infty=\Theta_1/3$ and
\begin{equation}
\Lambda_{\rm eff}=\frac{3H_\infty^2}{c^2}.
\end{equation}
For an asymptotic phase-dominated regime, or for the total effective fluid after
the usual Friedmann reduction, one may write
\begin{equation}
w_{\rm eff}(a)=-1-\frac{2}{3}\frac{\dot H}{H^2}.
\label{eq:weff_general}
\end{equation}
In a mixed matter-plus-dark-energy universe this is not automatically the
separate dark-energy equation of state unless matter, radiation, and curvature
have first been subtracted.
If the time derivative in Eq.~\eqref{eq:phaselaw} is the same cosmic-time
derivative, then
\begin{equation}
\dot H=\frac{1}{3}\ddot{\Theta},
\qquad
H^2=\frac{1}{9}\dot{\Theta}^{2},
\end{equation}
and hence
\begin{equation}
\boxed{w_{\rm eff}=-1-2\frac{\ddot{\Theta}}{\dot{\Theta}^{2}}.}
\label{eq:weff_cmc}
\end{equation}
Slow phase-rate drift gives $w_{\rm eff}\simeq-1$.

\begin{caveatbox}[Phase-law status]
The phase law is an ansatz, not a result derived in this paper.  A mathematical
completion must define the family of Hilbert spaces $\mathcal H_t$, the connection
$\nabla_t$, the phase $\Theta(t)$, the TT-sector state or density matrix
$\rho_{\rm TT}$, the regularization, and the time parameter.  The York-time
literature supports the use of York time in canonical cosmology, but does not by
itself establish the determinant-line phase law used here.  Earlier sign problems
in negative effective occupation numbers are therefore not repaired here; they
define the next calculation.
\end{caveatbox}

\section{Conservative Hybrid Parameterization}
\label{sec:hybrid}

The conservative model separates scale from evolution:
\begin{equation}
\rho_{\rm DE,eff}(a)=\rho_{\rm scale}F_{\rm dyn}(a),
\qquad
\rho_{\rm scale}=\rho_{\rm EH}^{(3)}(\Reff)=\rhoL,
\qquad
F_{\rm dyn}(1)=1.
\label{eq:hybrid}
\end{equation}
A minimal phenomenological proxy is
\begin{equation}
w(a)=-1+\varepsilon a^p,
\qquad
F_{\rm dyn}(a)=\exp\left[-\frac{3\varepsilon}{p}(a^p-1)\right],
\qquad \varepsilon>0,\quad p>0.
\label{eq:proxy}
\end{equation}
At $a=1$, $w_0=-1+\varepsilon$ and $w_a=-\varepsilon p$ in the usual
$w(a)=w_0+w_a(1-a)$ convention.  The minus sign in $F_{\rm dyn}$ follows
from $d\rho/d\ln a=-3(1+w)\rho$; for $\varepsilon>0$ and $p>0$ the component
is quintessence-like, with larger density in the past than at $a=1$.

The closed-geometry Friedmann equation should use a separate global curvature
radius $R_c$:
\begin{equation}
H^2(a)=\frac{8\pi G}{3c^2}\left[\rho_{m0}a^{-3}+\rho_{r0}a^{-4}+\rhoL F_{\rm dyn}(a)\right]
-\frac{kc^2}{a^2R_c^2},
\qquad k=+1.
\label{eq:friedmann_corrected}
\end{equation}
The normalization radius $\Reff$ in Eq.~\eqref{eq:hybrid} is not $R_c$.

\begin{figure}[htbp]
\centering
\begin{subfigure}[b]{0.48\textwidth}
\includegraphics[width=\textwidth]{figures/fig6_eos_proxy.pdf}
\caption{Toy $w(a)$ proxy.}
\end{subfigure}
\hfill
\begin{subfigure}[b]{0.48\textwidth}
\includegraphics[width=\textwidth]{figures/fig7_Ez_residuals.pdf}
\caption{Expansion residuals.}
\end{subfigure}
\caption{Phenomenological diagnostics.  These panels illustrate the difference
between an excluded pure $w=-1/3$ curvature-fluid interpretation and a near-$\Lambda$
phase proxy.  They are not likelihood fits.}
\label{fig:eos}
\end{figure}

\section{Numerical Tables}
\label{sec:numerical}

\begin{table}[htbp]
\centering
\caption{Constants and derived targets used in the numerical checks.}
\label{tab:const}
\small
\begin{tabular}{@{}llll@{}}
\toprule
Quantity & Symbol & Value & Units/status \\
\midrule
Speed of light & $c$ & $2.9979245800\times10^8$ & exact SI \\
Reduced Planck constant & $\hbar$ & $1.0545718170\times10^{-34}$ & J s \\
Newton constant & $G$ & $6.6743000000\times10^{-11}$ & SI \\
Planck length & $\ellP$ & $1.6162550239\times10^{-35}$ & m \\
Planck energy density & $\rhoP$ & $4.6329467907\times10^{113}$ & J m$^{-3}$ \\
Golden ratio & $\vph$ & $1.6180339887$ & exact \\
Hubble constant & $H_0$ & $2.1842852411\times10^{-18}$ & s$^{-1}$ \\
Dark-energy fraction & $\OmL$ & $0.6850000000$ & baseline input \\
Hubble radius & $\RH$ & $1.3724968349\times10^{26}$ & m \\
Critical density & $\rhoC$ & $7.6689477678\times10^{-10}$ & J m$^{-3}$ \\
Observed DE density & $\rhoL$ & $5.2532292210\times10^{-10}$ & J m$^{-3}$ \\
$\CEH=3/(8\pi)$ & $\CEH$ & $0.1193662073$ & derived \\
Effective radius & $\Reff$ & $1.6583113220\times10^{26}$ & m \\
Observable angle & $\chi_H$ & $0.8276472679$ & rad \\
Ball fraction & $f_{\rm obs}$ & $0.1048612055$ & diagnostic \\
\bottomrule
\end{tabular}
\end{table}

\begin{table}[htbp]
\centering
\caption{Coefficient candidates in $\rho=Cc^4/(GR^2)$.}
\label{tab:coeff}
\small
\begin{tabular}{@{}lrrrl@{}}
\toprule
$C$ expression & $C$ & $R(C)$ [Gly] & $R(C)/\RH$ & Status \\
\midrule
$3/(8\pi)$ & 0.119366 & 17.53 & 1.208 & spatial EH coefficient \\
$3\OmL/(8\pi)$ & 0.081766 & 14.51 & 1.000 & observational coefficient \\
$1/(2\pi^2)$ & 0.050661 & 11.42 & 0.787 & geometric candidate only \\
$3/(4\pi^2)$ & 0.075991 & 13.99 & 0.964 & geometric candidate only \\
$1$ & 1.000000 & 50.73 & 3.497 & raw $\beta=2$ Planck scaling \\
\bottomrule
\end{tabular}
\end{table}

\begin{figure}[htbp]
\centering
\includegraphics[width=0.88\textwidth]{figures/fig8_Rstar_beta.pdf}
\caption{Matching radius as a function of $\beta$.  The integer branch $\beta=2$
lands in the cosmological-radius band; neighboring integer branches do not.}
\label{fig:Rstar}
\end{figure}

\section{Observational Protocol}
\label{sec:obs}

For the global closed geometry, use the standard relations
\begin{align}
D_C(z)&=c\int_0^z\frac{dz'}{H(z')},\\
D_M(z)&=R_c\sin\left(\frac{D_C(z)}{R_c}\right),\\
D_L(z)&=(1+z)D_M(z),\qquad D_A(z)=\frac{D_M(z)}{1+z}.
\end{align}
The model must be compared to supernova, BAO, and CMB data using these distances,
not by fitting $D_\gamma$ directly.  Pantheon+ provides 1701 light curves of
1550 distinct Type Ia supernovae across $z=0.001$--$2.26$~\cite{PantheonPlus2022}.
DESI DR2 BAO uses more than 14 million galaxies and quasars and, in combined
$w_0w_a$ fits, reports a favored quadrant $w_0>-1$, $w_a<0$ while still finding
that BAO distances are well described by flat $\Lambda$CDM~\cite{DESIDR2BAO2025}.

\begin{falsebox}[Falsifiability tests]
\begin{enumerate}[leftmargin=1.5em,itemsep=2pt]
\item \textbf{Equation-of-state sign and size.}  The proxy in Eq.~\eqref{eq:proxy}
with $\varepsilon>0$ and $p>0$ predicts $w_0>-1$ and $w_a<0$.  Strong evidence for the
opposite sign would falsify that realization.
\item \textbf{Pure curvature-fluid rejection.}  If the model is forced to evolve
as $\rho\propto a^{-2}$, it predicts $w=-1/3$ and must fail late-time acceleration
fits.
\item \textbf{Curvature bound.}  The global $R_c$ must satisfy CMB+BAO curvature
bounds.  $R_c=\Reff$ is excluded.
\item \textbf{Topology tests.}  Matched circles apply only to multiply connected
quotients $S^3/\Gamma$, not to simply connected $S^3$.  Non-detection constrains
quotient models, not the minimal simply connected model by itself~\cite{Cornish1998}.
\item \textbf{Low-mode spectrum.}  Compact positive curvature gives a discrete
Laplacian spectrum.  Any claimed low-$\ell$ CMB signature must be computed with
proper transfer functions and likelihoods.
\end{enumerate}
\end{falsebox}

\section{Internal Consistency Conditions}
\label{sec:corrections}

\begin{enumerate}[leftmargin=1.5em,itemsep=3pt]
\item \textbf{Volume normalization.}  Use $V_{S^3}=2\pi^2R^3$, not the Euclidean
ball volume $4\pi R^3/3$, except in the local $\chi\ll1$ limit.
\item \textbf{Coefficient separation.}  $\CEH=3/(8\pi)$ is geometrically derived.
$\Cobs=3\OmL/(8\pi)$ is observational until $\OmL$ is derived.
\item \textbf{Branch labels.}  The raw $\beta=2$ Planck scaling has coefficient
$C=1$ and equality radius $50.73\,\Gly$; the EH-normalized equality radius is
$17.53\,\Gly$.
\item \textbf{Radius separation.}  $\Reff$ must not be identified with the global
FLRW curvature radius $R_c$.
\item \textbf{85-level disambiguation.}  $n=85$ is not a geometric cosmic level.
It can only be an amplitude/calibration index if Eq.~\eqref{eq:ncal} is justified.
\item \textbf{Dimensional hygiene.}  Expressions such as $3R_n+H$ are invalid.
Use $c/R_n$, $H R_n/c$, or $3/R_n$ depending on the intended dimension.
\item \textbf{Matched circles.}  They are not a direct test of simply connected
$S^3$; they test multiply connected quotient identifications.
\item \textbf{CMC phase law.}  It is a candidate closure, not a completed derivation.
\end{enumerate}

\section{Critical Weak Points}
\label{sec:weak}

\begin{openpbox}[P1: Missing derivation of $\OmL$]
The ratio $\Cobs/\CEH=\OmL$ is inserted from observation.  A first-principles
compact-$S^3$ derivation is still absent.
\end{openpbox}

\begin{openpbox}[P2: Effective radius vs. global curvature radius]
The model survives curvature bounds only if $\Reff$ is an effective
curvature-energy scale, not the global FLRW curvature radius.  This separation
needs an action-level derivation.
\end{openpbox}

\begin{openpbox}[P3: No action for $F_{\rm dyn}$]
The hybrid factor $F_{\rm dyn}(a)$ is phenomenological until it is derived from a
Lagrangian, Hamiltonian constraint, or controlled semiclassical approximation.
\end{openpbox}

\begin{openpbox}[P4: Null-domain distances not fitted]
The identity for $D_\gamma$ is a coordinate-domain identity until it is mapped to
$D_C,D_M,D_L,D_A$, BAO distances, and redshift.  It is not yet a falsifiability
test by itself.
\end{openpbox}

\begin{openpbox}[P5: CMC phase computation incomplete]
A viable phase completion must specify $(\mathcal H_t,\nabla_t,\Theta(t),
\rho_{\rm TT})$, the regularization and time parameter, then compute
$\Theta(t)$ from a concrete compact quantum state while verifying positivity,
gauge invariance, and time-convention stability.
\end{openpbox}

\section{Final Verdict}
\label{sec:verdict}

\begin{tcolorbox}[colback=green!4,colframe=falsecol,fonttitle=\bfseries\large,
  title=Four-Level Verdict,boxrule=1.0pt,arc=4pt]
\textbf{Level 1 -- established algebraic scaling:}
\[
\rho_P\ell_P^2=\frac{c^4}{G},
\qquad
\beta=2 \quad\Longleftrightarrow\quad
\rho_2(R)=\frac{c^4}{G R^2}.
\]
Among the tested integer branches $\beta=1,2,3,4$, this is the only one that
reaches the observed dark-energy density at cosmological radii.

\medskip
\textbf{Level 2 -- established scalar-curvature normalization:}
\[
\int_{S^3}\sqrt{g}\,{}^{(3)}R\,d^3x=12\pi^2R
\quad\Longrightarrow\quad
\rho_{\rm EH}^{(3)}(R)=\frac{3}{8\pi}\frac{c^4}{GR^2}.
\]
This is a positive constant-curvature scalar-density scale, not the negative
effective curvature fluid of the closed Friedmann equation, and not a uniquely
topological density coefficient.

\medskip
\textbf{Level 3 -- observationally inserted coincidence:}
\[
\rho_{\rm EH}^{(3)}(\Reff)=\rho_\Lambda,
\qquad
\Reff=\frac{\RH}{\sqrt{\OmL}}=17.53\,\Gly.
\]
This identity is exact after inserting $\OmL$: $\Reff=c/(H_0\sqrt{\OmL})=c/H_\Lambda$
is the observed dark-energy length in this normalization.  Thus the equality is a
curvature-scale coincidence, not a derivation from compact geometry.

\medskip
\textbf{Level 4 -- open dynamics:}
A conserved $R^{-2}$ component has $w=-1/3$ and does not drive late-time
acceleration.  The remaining problem is dynamical: derive $F_{\rm dyn}(a)$ from
an action, phase law, or null-domain likelihood construction, or reject the
model.
\end{tcolorbox}

\begin{quote}
\itshape
The paper establishes a robust curvature-scale coincidence: the unique tested
integer density branch $\beta=2$ gives $\rho\propto R^{-2}$, and the
constant-positive-curvature scalar-density calculation fixes the natural coefficient
$3/(8\pi)$.  Compactness makes the total integral finite, but the local density
coefficient is not uniquely topological.  Evaluated at the observationally defined
radius $\RH/\sqrt{\OmL}=c/H_\Lambda$, this reproduces $\rho_\Lambda$ because that
radius is defined from the observed dark-energy scale.  This is not a derivation
of dark energy, because $\OmL$ and the effective radius are inserted from the
Friedmann decomposition, and a conserved $R^{-2}$ component has $w=-1/3$.  The
remaining problem is dynamical: derive $F_{\rm dyn}(a)$, or reject the model.
\end{quote}

\section{Next Calculations Required}
\label{sec:next}

\begin{enumerate}[leftmargin=1.5em,itemsep=3pt]
\item \textbf{ADM closure.}  Derive or reject Eq.~\eqref{eq:rho_EH} inside the
full ADM Hamiltonian constraint with lapse, extrinsic curvature, and reference terms.
\item \textbf{$\OmL$ derivation.}  Test whether $\OmL$ can arise from a phase-sector
ratio, renormalized curvature residue, or observable-domain mapping without
circular input.
\item \textbf{Null-domain likelihood.}  Map $D_\gamma$ to $D_C,D_M,D_L,D_A$ and fit
Pantheon+ and DESI DR2.
\item \textbf{CMC phase calculation.}  Compute $\Theta(t)$ for a concrete compact
$S^3$ TT bundle or density matrix and verify Eq.~\eqref{eq:phaselaw}.
\item \textbf{Closed-mode Boltzmann implementation.}  Implement $S^3$ mode sums in
CLASS/CAMB-like pipelines and compare low-$\ell$ spectra to Planck likelihoods.
\item \textbf{Model comparison.}  Run Bayesian comparison against $\Lambda$CDM for
parameters $(\varepsilon,p,R_c)$ with conservative priors, including $\varepsilon>0$ and $p>0$ for the proxy tested here.
\end{enumerate}

\appendix

\section{Notation Dictionary}
\label{app:notation}

\begin{center}
\small
\begin{tabular}{@{}lll@{}}
\toprule
Symbol & Meaning & First defined \\
\midrule
$\vph$ & golden ratio $(1+\sqrt5)/2$ & Sec.~\ref{sec:hier} \\
$R_n=\ellP\vph^n$ & $\vph$-hierarchy radius & Sec.~\ref{sec:hier} \\
$n_R$ & geometric level $\ln(R/\ellP)/\ln\vph$ & Sec.~\ref{sec:hier} \\
$n_A=n_R-207$ & amplitude-level hypothesis & Sec.~\ref{sec:hier} \\
$\CEH=3/(8\pi)$ & spatial EH curvature coefficient & Sec.~\ref{sec:EH} \\
$\Reff=\RH/\sqrt{\OmL}$ & effective curvature-energy radius & Sec.~\ref{sec:tautology} \\
$R_c$ & global FLRW curvature radius & Sec.~\ref{sec:radii} \\
$\chi_H=\sqrt{\OmL}$ & Hubble angle on auxiliary $S^3(\Reff)$ & Sec.~\ref{sec:radii} \\
$K(\chi)$ & $S^3$ volume-growth kernel & Sec.~\ref{sec:geom} \\
$D_\gamma$ & light-front arc-length variable & Sec.~\ref{sec:nulldom} \\
$F_{\rm dyn}(a)$ & phenomenological DE evolution factor & Sec.~\ref{sec:hybrid} \\
\bottomrule
\end{tabular}
\end{center}

\section{Independent Numerical Check}
\label{app:numeric}

The central identities can be reproduced from
\begin{align}
\ellP&=\sqrt{\frac{\hbar G}{c^3}},
&\rhoP&=\frac{c^7}{\hbar G^2},
&\rhoP\ellP^2&=\frac{c^4}{G},\\
\rhoC&=\frac{3H_0^2c^2}{8\pi G},
&\rhoL&=\OmL\rhoC,
&\Reff&=\sqrt{\frac{\CEH c^4}{G\rhoL}}.
\end{align}
With the values of Table~\ref{tab:const},
\begin{equation}
\frac{\rho_{\rm EH}^{(3)}(\Reff)}{\rhoL}=1.0000000000
\end{equation}
within rounding precision.

\begin{thebibliography}{9}

\bibitem{Planck2018}
N.~Aghanim \textit{et al.} (Planck Collaboration),
``Planck 2018 results. VI. Cosmological parameters,''
\emph{Astron. Astrophys.} \textbf{641}, A6 (2020), arXiv:1807.06209.

\bibitem{Wald1984}
R.~M. Wald,
\emph{General Relativity}, University of Chicago Press, Chicago (1984).

\bibitem{Perelman2002}
G.~Perelman,
``The entropy formula for the Ricci flow and its geometric applications,''
arXiv:math/0211159 (2002).

\bibitem{Perelman2003}
G.~Perelman,
``Ricci flow with surgery on three-manifolds,''
arXiv:math/0303109 (2003).

\bibitem{RoserValentini2014}
P.~Roser and A.~Valentini,
``Classical and quantum cosmology with York time,''
\emph{Class. Quantum Grav.} \textbf{31}, 245001 (2014).


\bibitem{DESIDR2BAO2025}
M.~Abdul Karim \textit{et al.} (DESI Collaboration),
``DESI DR2 Results II: Measurements of Baryon Acoustic Oscillations and Cosmological Constraints,''
\emph{Phys. Rev. D} \textbf{112}, 083515 (2025), arXiv:2503.14738.

\bibitem{PantheonPlus2022}
D.~Brout \textit{et al.},
``The Pantheon+ Analysis: Cosmological Constraints,''
\emph{Astrophys. J.} \textbf{938}, 110 (2022), arXiv:2202.04077.

\bibitem{Cornish1998}
N.~J. Cornish, D.~N. Spergel, and G.~D. Starkman,
``Circles in the sky: finding topology with the microwave background radiation,''
\emph{Class. Quantum Grav.} \textbf{15}, 2657--2670 (1998), arXiv:astro-ph/9801212.


\end{thebibliography}

\end{document}
